%% kapa: exam=1 round=28 year=2017 period=2 form=A subject=appraisal_law subject_ko=감정평가관계법규 seq=9 number=9
%% source: Q-Net 감정평가사 기출문제 (cst00309) artlSeq=5199603 file=extracted/1008/5199603/2017년 제28회 감정평가사 1차시험 문제지/2교시 A형.pdf
%% conversion: kordoc 4.24.51 -> tools/gen1.py (auto); review=proofread-sonnet; flags=none
\begin{question}{9}
국토의 계획 및 이용에 관한 법령상 입지규제최소구역 등에 관한 설명으로 옳지 않은 것은?
\begin{choices}
  \item 시ㆍ도지사는 도시ㆍ군기본계획에 따른 도심ㆍ부도심 또는 생활권의 중심지역을 입지규제최소구역으로 지정할 수 있다.
  \item 입지규제최소구역계획에는 입지규제최소구역의 지정 목적을 이루기 위하여 건축물의 건폐율ㆍ용적률ㆍ높이에 관한 사항이 포함되어야 한다.
  \item 입지규제최소구역계획은 도시ㆍ군관리계획으로 결정한다.
  \item 입지규제최소구역에 대하여는「주택법」제35조에 따른 대지조성기준을 적용하지 아니할 수 있다.
  \item 입지규제최소구역으로 지정된 지역은「건축법」에 따른 특별건축구역으로 지정된 것으로 본다.
\end{choices}
\end{question}
\answer{1}
